\documentclass[letterpaper,11pt,leqno]{article}
\usepackage{paper}

\hypersetup{pdftitle={Research Statement -- Mervin Goklas Hamonangan}}

\newcommand{\bib}{research-statement.bib}

\begin{document}

\title{Research Statement}

\author{Mervin Goklas Hamonangan
\thanks{Institute for Economic and Social Research (LPEM), Faculty of Economics and Business, Universitas Indonesia. \url{mervin.goklas@gmail.com}}}

\date{September 2026}

\paperurl{https://mervin-27.github.io/}

\begin{titlepage}
\maketitle

I work at the intersection of political economy and the economics of religion, using historical and administrative data to study how institutions shape religiosity and religious-political alignment, and how those alignments in turn shape development outcomes. My graduate research examines how colonial education policy in the Dutch East Indies shaped the geography of religious and communist political support after independence. A research assistantship in economic history and a governance-focused consulting engagement with the World Bank keep this agenda grounded in how institutional data of this kind is built and used in practice.

\end{titlepage}

\section{Research agenda}\label{s:agenda}

My work is organized around one question: how do the institutions a state inherits, builds, or is forced to build determine the ideological alignments of its population, and how do those alignments feed back into the policies and outcomes available to it. Two literatures inform this question. The first traces the long-run political effects of colonial institutions (Dell and Olken, 2020; Cantoni et al., 2017). The second traces the effect of education on religiosity and political preference (Gulesci and Meyersson, 2013; Hungerman, 2014; Cesur and Mocan, 2017). Both are well developed for Europe, the Americas, and parts of South Asia. Neither is well developed for maritime Southeast Asia, where a mixed colonial legacy, an unusually plural religious landscape, and one of the most consequential ideological ruptures of the twentieth century, the 1965--66 anti-communist violence in Indonesia, sit on top of each other. My aim is to fill that gap with data that has not previously been assembled at this resolution.

\section{Colonial education and post-independence political ideology}\label{s:education}

My LSE master's extended essay, \textit{Colonial Education Reforms and Political Ideology post-Independence: The Ethical Policy of Dutch East Indies}, is the first piece of this agenda. The Dutch colonial government's 1901 Ethical Policy expanded both elite European-language schooling and mass village schooling across the East Indies, ostensibly to return some of the wealth extracted under the earlier Cultivation System (\textit{Cultuurstelsel}) to the indigenous population (van Zanden, 2010; Ricklefs, 2008). The policy is a well-known example of how colonial nation-building can fail on its own terms: it produced a small, educated elite that went on to lead the anticolonial movement (van Niel, 1970).

I ask whether the intensity of this school expansion left a distinguishable ideological imprint on the districts that received it, using Indonesia's first and most competitive democratic election, held in 1955, as the outcome. I combine three novel-to-this-question datasets: district-level schooling and literacy counts from the 1930 \textit{Volkstelling}, the last and largest Dutch colonial census; district-level vote shares for the Islamist Masyumi party, Nahdlatul Ulama, and the Communist Party of Indonesia (PKI) from the 1955 election, shared by Samuel Bazzi; and 1920 pesantren, mosque, and Darul Islam rebellion data assembled by Bazzi et al. (2020). Using standard inference on this merged panel, I find that districts with greater colonial-era school exposure show significantly lower support for religious-based parties in 1955, a result consistent with Cantoni et al.'s (2017) finding that curriculum content shapes political ideology through channels beyond simple literacy. The essay received a mark of Distinction.

Two extensions follow directly. The first pushes past the reduced-form vote-share result to isolate mechanism: does school exposure operate through literacy, through the language of instruction, or through exposure to a colonial bureaucratic elite track (OSVIA, STOVIA) that created cross-cutting loyalties independent of religious networks. The second brings the outcome forward from 1955 to the 1965--66 violence itself, asking whether the same colonial-era schooling variation predicts where anti-communist mobilization was most and least severe, a question that speaks directly to the debate over local versus state-driven explanations of the killings (Roosa, 2006).

\section{Religiosity, education, and the secularization hypothesis}\label{s:religion}

A second, related strand sits inside the broader economics of religion literature (Iyer, 2016). The secularization hypothesis holds that economic development, through education and urbanization, reduces religiosity; the evidence is mixed and mostly Western. Gulesci and Meyersson (2013) find that a compulsory-schooling reform in Turkey reduced women's religiosity; Hungerman (2014) finds the same pattern for Canada; but Squicciarini (2020) finds that religiosity slowed the adoption of technical education and growth in nineteenth-century France, the reverse mechanism. My extended essay's finding, that colonial school exposure predicts lower support for religious-based parties in Indonesia, sits on the secularization side of this debate, but the mechanism is unresolved: whether education reduced religiosity itself, or only its translation into party-political support, remains an open question. I intend to take on this question directly with contemporary Indonesian survey and voting data, extending the analysis from the 1955 election into the post-Reformasi period.

\section{Research experience}\label{s:experience}

Three engagements shape how I approach this agenda. Since July 2025, I have been a research assistant in economic history at the Faculty of Economics and Business, Universitas Indonesia, working with historical sources of the same kind as the colonial census and election returns behind the extended essay. From July to December 2024, I was a short-term consultant with the World Bank's Governance Global Practice, where I saw how institutional and governance measures are constructed, and where they break down, before being used in policy analysis. Earlier, as a policy and research intern at J-PAL Southeast Asia (2021) and a research assistant at LPEM FEB UI (2022--2023), I worked on program evaluation and survey data, which is where I learned to ask what a given identification strategy can and cannot rule out. Together they explain why this agenda leans on archival data and careful measurement.

\section{Next steps}\label{s:next}

My immediate priority is to convert the master's essay into a standalone paper, extending the sample and sharpening the identification strategy around specific school-type openings rather than aggregate exposure. In parallel, I am scoping the contemporary religiosity data needed to test the secularization mechanism directly. Both projects are designed to travel: the colonial-education question generalizes to other former colonies with comparable census and election archives, and the secularization question generalizes to any setting where religious and political affiliation can be jointly observed over time. I am looking to pursue this agenda in a PhD program with strength in political economy, economic history, and the economics of religion.

\end{document}
